\maketitle

\begin{abstract}
A thematic survey of the literature bearing on the question:
\emph{when does observationally coherent data determine a genuine
probability measure, and what additional commitments does this
require?}  The review is organized around four themes: (1)~the
boundary between finite additivity and $\sigma$-additivity,
(2)~Stone duality and the problem of realization,
(3)~value-definiteness and obstruction in non-distributive algebras,
and (4)~delay reconstruction as a source of observation algebras.
Within each theme, classical results are placed alongside the
philosophical, logical, and topos-theoretic perspectives that
illuminate the boundary.  Open problems and closed directions are
documented where they arise.  This is a working survey, not a
journal submission; it occasionally marks my own contributions and
dead ends alongside the literature.
\end{abstract}

\tableofcontents

% ==================================================================
\section{Introduction}
\label{sec:intro}
% ==================================================================

The central question of this survey decomposes into three layers:
\begin{enumerate}
  \item \textbf{Extension:} When do compatible finitely additive
    charges on a directed system of Boolean algebras determine a
    $\sigma$-additive probability measure?
  \item \textbf{Realization:} When does a measure on a dual space
    (Stone space) descend to a measure on an underlying state space?
  \item \textbf{Value-definiteness:} When does a probabilistic
    description admit simultaneous definite values for all
    observables?
\end{enumerate}
Each layer requires a commitment the previous does not.  The
algebraic structure of the observation algebra constrains which
commitments are available; distributivity is sufficient to make all
three available.

The survey draws on measure theory, Boolean algebra, Stone duality,
orthomodular lattice theory, topos theory, philosophy of
probability, model theory, and dynamical systems.  Coverage is not
exhaustive within any one field; the aim is to document the results
that bear on the three-layer decomposition, including the negative
results that close off natural research directions.

\medskip\noindent
\textbf{Scope.}  The review covers:
\begin{itemize}[nosep]
  \item The classical extension theorems (Carath\'{e}odory,
    Kolmogorov, Choksi) and their algebraic characterization
    (Yosida--Hewitt, Kantorovich--Vulikh--Pinsker);
  \item The model-theoretic obstruction to first-order
    characterization of $\sigma$-additivity
    (\L{}o\'{s}, Hoover, Keisler);
  \item The philosophical context (de Finetti, Howson, Cox) and the
    logical coordinates (Frot, Esp\'{i}ndola, Caramello);
  \item Stone duality and the realization problem
    (Stone, Fremlin, Rao--Rao);
  \item Orthomodular lattice theory and contextuality
    (Kochen--Specker, Gleason, Pt\'{a}k--Pulmannov\'{a},
    McDonald--Bimb\'{o});
  \item Delay embedding and reconstruction
    (Takens, Krieger, Sauer--Yorke--Casdagli);
  \item Sheaf-theoretic and topos-theoretic approaches to the
    extension problem (Zafiris, Biesel, Caramello, Vickers);
  \item Four open problems at the boundary.
\end{itemize}

% ==================================================================
\section{The extension boundary: charges to measures}
\label{sec:extension}
% ==================================================================

\subsection{Classical extension theorems}

The foundational results on extending finitely additive charges to
$\sigma$-additive measures are well established:

\begin{itemize}
  \item \textbf{Continuity from above.}  A finitely additive charge
    $\ell$ on an algebra $\mathcal{A}$ is $\sigma$-additive if and
    only if $E_n \downarrow \varnothing$ implies
    $\ell(E_n) \to 0$~\cite[\S13, Theorem~D]{Halmos1950}.

  \item \textbf{Carath\'{e}odory extension.}  A $\sigma$-additive
    premeasure on a semiring extends uniquely to a measure on the
    generated $\sigma$-algebra~\cite[\S13]{Halmos1950},
    \cite[Vol.~1, Ch.~1]{Bogachev2007}.

  \item \textbf{Kolmogorov extension.}  Compatible
    finite-dimensional distributions on a product space extend to a
    unique probability measure under standard Borel/tightness
    conditions~\cite[Ch.~36]{Billingsley1995}.

  \item \textbf{Choksi inverse limits.}  Compatible measures on an
    inverse system of compact Hausdorff spaces extend to the inverse
    limit~\cite[Theorem~1]{Choksi1958}.
\end{itemize}

The \textbf{Yosida--Hewitt decomposition}~\cite{YosidaHewitt1952}
gives the structural anatomy: every bounded finitely additive charge
decomposes uniquely as $\ell = \ell_c + \ell_p$, where $\ell_c$ is
$\sigma$-additive and $\ell_p$ is purely finitely additive (PFA).
The extension problem reduces to: when is $\ell_p = 0$?

\subsection{The algebraic characterization}
\label{sec:kvp}

The deepest algebraic answer is
Fremlin's~\cite[\S391]{Fremlin2003} packaging of results going back
to Kelley and to the theory of weakly
$(\sigma,\infty)$-distributive algebras: a Boolean algebra
$\mathfrak{A}$
admits a strictly positive $\sigma$-additive probability if and only
if $\mathfrak{A}$ is
\begin{enumerate}[label=(\roman*),nosep]
  \item Dedekind $\sigma$-complete,
  \item weakly $(\sigma,\infty)$-distributive, and
  \item chargeable (admits a strictly positive finitely additive
    functional).
\end{enumerate}
Kelley's theorem~\cite[391J]{Fremlin2003} characterizes
chargeability combinatorially.

Condition~(ii) is the algebraic face of $\sigma$-additivity: it is
a genuine lattice condition on how countable infima interact with
arbitrary suprema, not a tautological restatement.  This is the
structural answer to the question of what property of the
\emph{algebra} (the site) --- rather than the charges (the
sections) --- controls whether local data extend to a
$\sigma$-additive global measure.

The boundary between chargeability and measurability was sharpened by
Talagrand~\cite{Talagrand2008}, who constructed a Maharam submeasure
that is not equivalent to any measure, resolving Maharam's 1947
control problem.  This shows that the gap between conditions~(iii)
and~(i)--(iii) together is genuine: an algebra can be chargeable
without admitting any $\sigma$-additive probability.

\subsection{Non-axiomatizability}

A finitely additive charge on a Boolean algebra satisfies all
first-order conditions that each individual measure in a family
satisfies, by \L{}o\'{s}'s theorem~\cite{Los1955}.  Applied to
Dirac masses $(\delta_n)_{n \in \N}$ on $\mathcal{P}(\N)$ with a
nonprincipal ultrafilter $\mathcal{U}$: each $\delta_n$ is
$\sigma$-additive, but the ultraproduct charge $\ell(A) =
\lim_\mathcal{U}\delta_n(A) = \mathbf{1}_{A \in \mathcal{U}}$ is
the $\{0,1\}$-valued ultrafilter charge, which is purely finitely
additive.  Hence no first-order theory of Boolean algebras with
charge can characterize $\sigma$-additivity
(Hoover~\cite{Hoover1978}; Fagin--Halpern--Megiddo~\cite{FaginHalpernMegiddo1990};
see also Keisler~\cite{Keisler1985}).

The same technique establishes non-axiomatizability at the level of
pure Boolean algebra language (without charge predicates): the
Lebesgue measure algebra and the Cohen algebra
$\mathrm{RO}(2^{<\omega})$ are both atomless, hence elementarily
equivalent by Tarski's classification, yet only the former admits a
strictly positive $\sigma$-additive probability.  (It satisfies all
three KVP conditions; the Cohen algebra fails weak
$(\sigma,\infty)$-distributivity.)  This folklore observation
appears implicit in Koppelberg~\cite[Ch.~18]{Koppelberg1989}.

\subsection{The philosophical context}
\label{sec:philosophical}

De Finetti~\cite{deFinetti1974} grounds probability in
\emph{coherence} (the Italian \emph{coerenza}, which translates
equally well as ``coherence'' or ``consistency''; the English
literature settled on ``coherence'').  His criterion ---
no finite combination of bets at the given rates guarantees a
loss --- characterizes exactly the finitely additive probability
functions.  De Finetti maintained that $\sigma$-additivity is a
mathematical convenience with no operational justification: the
Dutch Book argument for it requires extending the ``finite sum of
fair bets is fair'' postulate to countable sums, which circularly
entails $\sigma$-additivity~\cite[\S\S3--4]{Howson2008}.

Howson~\cite{Howson2008} makes de Finetti's position precise by
identifying two structural properties that de Finetti's consistency
shares with first-order logic:
\begin{itemize}[nosep]
  \item \emph{Absoluteness:} consistency of $P$ depends only on
    $P[E]$, not on which algebra $E$ sits in.
  \item \emph{Compactness:} inconsistency is detected by a finite
    subset of the assignments.
\end{itemize}
Both fail for $\sigma$-additivity.  Howson identifies a missing
``completeness theorem telling us that the rules of probability
extend to finite but not countable additivity''~\cite[p.~17]{Howson2008}.

Cox's axioms~\cite{Cox1946} produce finitely additive probability
from three scale-free structural principles plus propositional
logic.  Howson notes~\cite[\S6]{Howson2008}: ``Cox's result does not
extend to endorsing countable additivity.  There are infinitary
propositional languages\ldots\ $L_{\omega_1,\omega}$\ldots\ but
there is no natural extension of Cox's proof which would exploit
that additional structure.''  Ross~\cite{Ross2012} argues from the
opposite direction that ``all roads lead to violations of countable
additivity,'' reinforcing the view that $\sigma$-additivity requires
commitment beyond coherence.

\subsection{Logical coordinates}
\label{sec:logical}

Three results locate the extension boundary within formal logic and
topos theory.

\medskip\noindent
\textbf{Frot (2013).}  Deligne's theorem (``coherent topoi have
enough points'') is equivalent to G\"{o}del's completeness theorem
for first-order logic~\cite{Frot2013}.  The dictionary: models $=$
points of the classifying topos; consistency $=$ non-triviality;
coherent theory $=$ finitary axioms.  Finite additivity is a
coherent condition; $\sigma$-additivity is not (it requires
countable universal quantification).  So Deligne's theorem
guarantees the existence of finitely additive charges but is silent
about $\sigma$-additive measures.

\medskip\noindent
\textbf{Esp\'{i}ndola (2020).}  Generalizes Deligne to
$\kappa$-geometric logic~\cite{Espindola2020}: at
$\kappa = \omega_1$, $\sigma$-additivity becomes expressible.
Property~$T$ (an exactness condition combining strong
distributivity and dependent choice) is the structural content at
each level.  At $\kappa = \omega$, Property~$T$ is automatic,
recovering Deligne.  At $\kappa = \omega_1$, it is a genuine
restriction --- but applied to the site of finite Boolean
subalgebras where $\sigma$-additivity lives, Property~$T$ reduces
to $\sigma$-completeness of the underlying algebra plus uniqueness
of the Carath\'{e}odory extension.  Both are classical conditions
already assumed in any CE setup.  The machinery does not produce new
content at this specific site.

\medskip\noindent
\textbf{Caramello (2018).}  The bridge technique~\cite{Caramello2018}
transfers properties between syntactic and semantic sites via Morita
equivalences.  But $\sigma$-additivity cannot be expressed as a
geometric sequent (which permits only finite conjunction and small
disjunction), so it cannot appear as a quotient of any geometric
theory in Caramello's framework.

\subsection{Synthesis: five perspectives on the same boundary}
\label{sec:extension-synthesis}

A remarkable convergence emerges.  Five independent traditions ---
operating with different vocabularies, motivations, and proof
techniques --- identify the same structural boundary:

\begin{enumerate}
  \item \textbf{Philosophy of probability} (de Finetti, Howson):
    consistency (finite additivity) has absoluteness and compactness;
    $\sigma$-additivity has neither.  Cox's axioms produce finite
    additivity only.

  \item \textbf{Logic and topos theory} (Frot, Esp\'{i}ndola,
    Caramello): Deligne (coherent topoi have enough points)
    $=$ G\"{o}del completeness $=$ finitary completeness.
    $\sigma$-additivity escapes coherence.  The
    $\omega_1$-coherent level can express it, but the machinery
    collapses to classical conditions on the relevant site.

  \item \textbf{Algebraic measure theory} (Fremlin, KVP): weak
    $(\sigma,\infty)$-distributivity is the lattice condition that
    forces $\sigma$-additivity, answering Howson's question in
    algebraic language.

  \item \textbf{Model theory} (\L{}o\'{s}, Hoover, Keisler):
    ultraproducts kill first-order axiomatizability.  The class of
    $\sigma$-additive structures is not elementary.

  \item \textbf{Conglomerability and imprecise probability}
    (Seidenfeld et al.~\cite{Seidenfeld1984},
    Walley~\cite{Walley1991}): conglomerability --- the requirement
    that conditional probabilities given a partition constrain the
    unconditional probability --- is the operational face of
    $\sigma$-additivity.  Schervish--Seidenfeld--Kadane show its
    failure is pervasive under finite additivity; Walley develops
    the imprecise-probability framework where the distinction is
    central.
\end{enumerate}

No published work unifies these five perspectives under a single
theorem.  Whether such a unification would constitute genuine
mathematics or terminological packaging is an open question.  At
minimum, the convergence confirms that the finite/countable boundary
is a robust structural phenomenon, not an artifact of any particular
framework.

\subsection{Sheaf-theoretic approaches: closed directions}
\label{sec:sheaf-closed}

Several natural attempts to reformulate $\sigma$-additivity as a
sheaf or descent condition have been investigated and found to yield
dictionary translations rather than theorems:

\begin{itemize}
  \item \textbf{Zafiris (2006).}  The Grothendieck topology $J$
    (epimorphic families) on the site of Boolean subalgebras does
    not grade covers by cardinality~\cite{Zafiris2006}.  Finite,
    countable, and uncountable covers satisfy $J$ uniformly; the
    finite/countable boundary is invisible.

  \item \textbf{Biesel (2024).}  Measures form a sheaf for finite
    covers~\cite[Theorem~8]{Biesel2024}; probability measures
    likewise~\cite[Theorem~11]{Biesel2024}.  Countable covers are
    explicitly excluded~\cite[Remark~7]{Biesel2024}.  The step
    from finite to countable covers is exactly the
    $\sigma$-additivity boundary, but this is tautological:
    finitely additive $=$ sheaf for finite covers;
    $\sigma$-additive $=$ sheaf for countable covers.

  \item \textbf{Caramello (2018).}  The bridge technique cannot
    express $\sigma$-additivity (\S\ref{sec:logical}).

  \item \textbf{Vickers (2008).}  The localic theory of
    valuations~\cite{Vickers2008} provides the framework for
    formulating ``a valuation is supported on the sublocale of
    principal ultrafilters'' as a frame-internal condition, but
    this formulation has not been carried out.  The earlier
    textbook~\cite{Vickers1989} develops locales and sublocales
    without addressing measures.
\end{itemize}

The pattern is consistent: every finitary framework (sheaf-theoretic,
topos-theoretic, locale-theoretic) either cannot see the
finite/countable boundary or restates $\sigma$-additivity as a
definition.  The one framework that can express it (Esp\'{i}ndola's
$\omega_1$-geometric logic) collapses to classical conditions on
the relevant site.

\subsection{Open problems at the extension boundary}

\subsubsection{The $\sigma$-complete characterization pattern}

Every known characterization theorem for when a Boolean algebra
carries a strictly positive $\sigma$-additive probability
presupposes Dedekind $\sigma$-completeness:
\begin{itemize}[nosep]
  \item Kelley~\cite{Kelley1959}, Theorem~9: complete + weakly
    distributive + chargeable.
  \item KVP~\cite[391D]{Fremlin2003}: Dedekind $\sigma$-complete +
    weakly $(\sigma,\infty)$-distributive + chargeable.
  \item Balcar--Jech--Paz\'{a}k~\cite[539N]{Fremlin2003} (under
    p-ideal dichotomy): Dedekind complete + ccc $\Rightarrow$
    Maharam algebra.
  \item Kalton--Roberts~\cite[392F--G]{Fremlin2003}: uniformly
    exhaustive submeasure + Dedekind $\sigma$-complete.
  \item All negative constructions (Argyros~\cite{Argyros1983},
    Souslin line under $\neg$SH, Farah--Veli\v{c}kovi\'{c},
    Talagrand~\cite{Talagrand2008}) also operate in the
    $\sigma$-complete row.
\end{itemize}
The non-$\sigma$-complete regime has no analogous characterization ---
only piecewise existence results for specific construction families.
This makes Strategy~D (Problem~\ref{op:strategy-d} below)
structurally novel: not ``unstudied,'' but \emph{no
measure-existence characterization reaches the non-$\sigma$-complete
non-atomic regime}.  Confirmed across $\geq 6$ sources: Fremlin
ch39, ch53, ch54; Kelley~\cite{Kelley1959};
Argyros~\cite{Argyros1983};
D\v{z}amonja--Plebanek~\cite{DzamonjaPleb2008}.

\subsubsection{Strategy~D}

\begin{openproblem}[Measure-free Boolean algebra --- Strategy~D]
\label{op:strategy-d}
Does there exist a Boolean algebra $B$ satisfying:
\begin{enumerate}[label=(\roman*),nosep]
  \item $B$ is non-$\sigma$-complete
    $($equivalently, $\St(B)$ is not basically disconnected$)$;
  \item $B$ is non-atomic
    $($equivalently, $\St(B)$ has no isolated points$)$;
  \item $B$ is measure-free: $B$ admits no strictly positive
    $\sigma$-additive probability
    $($equivalently, $\St(B)$ carries no strictly positive Radon
    probability measure$)$?
\end{enumerate}
Here $\sigma$-additivity on a non-$\sigma$-complete $B$ is
interpreted for disjoint countable families whose join exists in $B$.
\end{openproblem}

On the clopen algebra of any Stone space, $\sigma$-additivity is
vacuous by compactness; the real question concerns Radon measures on
the Stone space itself.

\medskip\noindent
\textbf{Near-miss constructions.}  Table~\ref{tab:near-miss}
summarizes the known candidates and which Strategy~D conditions they
satisfy.

\begin{table}[ht]
\centering
\begin{tabular}{@{}lcccl@{}}
\toprule
\textbf{Construction} &
  \textbf{(i) atomless} & \textbf{(ii) non-$\sigma$-c.} &
  \textbf{(iii) measure-free} & \textbf{Notes} \\
\midrule
Argyros~\cite{Argyros1983} &
  \checkmark & \texttimes &
  \checkmark & under CH; $\beta Y$ route \\
Kunen L-space~\cite{Plebanek2023} &
  \checkmark & \checkmark &
  \texttimes & has s.p.\ measure by design \\
Souslin line (under $\neg$SH) &
  \checkmark & \texttimes &
  \checkmark & consistency-only \\
Gaifman algebra &
  \texttimes & \texttimes &
  \checkmark\textsuperscript{*} &
  \textsuperscript{*}no FA measure \\
Todor\v{c}evi\'{c}/D-P~\cite{DzamonjaPleb2008} &
  \texttimes & ? & \checkmark & ZFC; has atoms \\
\bottomrule
\end{tabular}
\caption{Near-miss constructions for Strategy~D.  Each fails at
  least one of the three required conditions.}
\label{tab:near-miss}
\end{table}

\noindent
Three structural patterns emerge:
\begin{enumerate}[nosep]
  \item The Stone--\v{C}ech compactification route
    (Argyros~\cite{Argyros1983}) naturally produces basically
    disconnected spaces, i.e., $\sigma$-complete Boolean algebras.
    This is a systematic obstruction: the standard technique for
    building measure-free spaces also forces $\sigma$-completeness.
  \item The L-space route (Kunen~\cite{Plebanek2023}) is
    consistency-only: under MA\,+\,$\neg$CH, no compact L-spaces
    exist.  Even where they do exist, the construction's purpose is
    to \emph{produce} a strictly positive measure.
  \item The $\mathcal{P}(\omega)/\mathrm{fin}$ subalgebra route
    (Todor\v{c}evi\'{c}~\cite{Todorcevic2000},
    D\v{z}amonja--Plebanek~\cite{DzamonjaPleb2008}) avoids
    $\sigma$-completeness but produces atoms (see below).
\end{enumerate}
Every known ZFC measure-free Boolean algebra is either
$\sigma$-complete or has atoms.

\medskip\noindent
\textbf{The Todor\v{c}evi\'{c} example.}
D\v{z}amonja--Plebanek~\cite[Theorem~3.1]{DzamonjaPleb2008} extract
from Todor\v{c}evi\'{c}'s Theorem~8.4~\cite{Todorcevic2000} a
subalgebra $\mathfrak{A}$ of $\mathcal{P}(T)/\mathrm{fin}$
generated by a family $\{T_a : a \in A\}$, where $A \subseteq Z$
(zero-density subsets of $\N$) is well-ordered under $\subseteq^*$
with order-type $\mathrm{add}(\mathcal{N})$, and
$T_a = \{(s,m) \in T : a \cap m \subseteq s\}$.  They prove:
$|\mathfrak{A}| = \mathrm{add}(\mathcal{N})$, $\mathfrak{A}$ is ccc
but not $\sigma$-centred, and $\mathfrak{A}$ carries no strictly
positive measure.  The last follows from Lemma~3.2: the generating
family satisfies $a \leq b$, $b \leq a$, or $a \cdot b = 0$ for
$a, b \in \mathcal{G}$; combined with non-$\sigma$-centredness, this
forces any strictly positive measure to assign positive mass to
finitely many maximal chains, a contradiction.

However, the generators $\{[T_a] : a \in A\}$ form a chain in
$\mathcal{P}(T)/\mathrm{fin}$: for $a \subseteq^* b$, replace $a$
by $a \cap b$ (same $\subseteq^*$-class, since $A$ is closed under
finite changes within $Z$) to get literal $a' \subseteq b$, whence
$T_b \subseteq T_{a'}$, so $[T_b] \leq [T_{a'}] = [T_a]$.  A
Boolean algebra generated by a single chain is an interval algebra.
Since $A$ has order-type $\mathrm{add}(\mathcal{N})$ (an ordinal),
immediate-successor pairs $a \prec a'$ produce atoms
$[T_a] \setminus [T_{a'}]$ in $\mathfrak{A}$.  Therefore
$\mathfrak{A}$ fails condition~(ii) --- it is not atomless.

\medskip\noindent
\textbf{Consistency results.}
D\v{z}amonja--Plebanek~\cite[Corollary~2.8]{DzamonjaPleb2008}
prove: under MA\,+\,$\neg$CH, every atomless ccc Boolean algebra of
size~$< \mathfrak{c}$ carries a strictly positive nonatomic measure.
Corollary~2.11 gives the converse: under MA\,+\,$\neg$CH, for
atomless algebras of size~$< \mathfrak{c}$, ccc is equivalent to
the chain condition of Theorem~2.9.  These results push Strategy~D
toward set-theoretic independence.

\medskip\noindent
\textbf{Current assessment.}
Plebanek's survey~\cite{Plebanek2024} confirms the exact parameter
regime (non-$\sigma$-complete, non-atomic, measure-free) remains
open.  Strategy~D is likely ZFC-independent, but this is not proved.
No ZFC construction or ZFC obstruction is known.%
\footnote{\textbf{Superseded 2026-06-18.} This regime in fact has a \emph{positive}
ZFC answer and was never field-open: Gaifman~1964 (PJM 14(1):61--73, Thm~2.2 $+$
property~$(\dagger)$) already exhibits an atomless Boolean algebra with no
strictly-positive \emph{finitely}-additive measure, in ZFC. The non-$\sigma$-complete
$+$ atomless conditions are trivial; the literature records these algebras only in
their \emph{completed} (extremally/basically disconnected) form, which is why the
regime ``looked'' open. ``Strategy~D'' was this programme's private name; retired.}

\subsubsection{Ultralimit representability}

\begin{openproblem}[Ultralimit representability --- reduces to
  Problem~\ref{op:strategy-d}]
\label{op:ultralimit}
Let $B$ be a Boolean algebra, $\mathcal{U}$ a nonprincipal
ultrafilter on an index set $I$, and $(\mu_i)_{i \in I}$ a family
of $\sigma$-additive probabilities on $B$.  Call a purely finitely
additive charge $\ell$ on $B$ \emph{ultralimit-representable} if
$\ell(b) = \lim_\mathcal{U}\mu_i(b)$ for all $b \in B$.

For which Boolean algebras $B$ is every purely finitely additive
charge ultralimit-representable?
\end{openproblem}

Settled in all regimes except Problem~\ref{op:strategy-d}:
\begin{itemize}[nosep]
  \item \emph{Atomic $B$} (e.g.\ $\mathcal{P}(\N)$): every PFA
    charge is representable (ultrafilter charges via Diracs; the
    rest by Krein--Milman).
  \item \emph{$\sigma$-complete $B$:} Nikod\'{y}m convergence
    \cite[IV.9.8]{DunfordSchwartz1958} forces ultralimits to stay
    $\sigma$-additive.  No PFA charge arises.
  \item \emph{Non-$\sigma$-complete, non-atomic $B$:} failure of
    representability is witnessed by a measure-free direct-product
    factor, reducing to Problem~\ref{op:strategy-d}.
\end{itemize}

\subsection{The operational boundary: estimation theory meets
  the finite/countable gap}
\label{sec:operational-boundary}

The finite/countable boundary is epistemically inaccessible from
finite data: no finitary test --- first-order sentence, finite
cover, sheaf condition on finite refinements --- separates
$\sigma$-additive measures from purely finitely additive charges
(\S\ref{sec:extension-synthesis}).  Model-theoretically, this is
Lindstr\"{o}m's theorem: first-order logic is the strongest logic
with compactness, and compactness is what ultraproducts exploit to
produce PFA charges from $\sigma$-additive ones.

Yet the boundary has operational consequences that are detectable
from finite samples.  The connection runs through minimax estimation
under contamination.

\medskip\noindent
\textbf{The Yosida--Hewitt / contamination correspondence.}
Huber's $\varepsilon$-contamination model assumes observations from
$(1-\varepsilon)f + \varepsilon g$, where $f$ is the target density
and $g$ is an unknown contamination.  This is structurally parallel
to the Yosida--Hewitt decomposition $\mu = \mu_c + \mu_p$: the
$\sigma$-additive part $\mu_c$ is the ``target'' and the PFA part
$\mu_p$ is the ``contamination'' --- a component that passes all
finite consistency checks but does not concentrate on any
$\sigma$-additive structure.

Liu--Gao~\cite{LiuGao2019} establish the minimax rate for density
estimation under contamination (Theorem~2.1):
\[
  \mathcal{R}(\varepsilon, \beta_0, \beta_1)
  \;\asymp\;
  \bigl[n^{-2\beta_0/(2\beta_0+1)}\bigr]
  \;\vee\;
  \bigl[\varepsilon^2(1 \wedge m)^2\bigr]
  \;\vee\;
  \bigl[n^{-2\beta_1/(2\beta_1+1)}\,
    \varepsilon^{2/(2\beta_1+1)}\bigr],
\]
where $\beta_0$ is the target smoothness, $\beta_1$ the
contamination smoothness, $\varepsilon$ the contamination
proportion, and $m$ the local contamination level.  Three terms:
\begin{enumerate}[nosep]
  \item $n^{-2\beta_0/(2\beta_0+1)}$: the classical minimax rate
    (no contamination);
  \item $\varepsilon^2(1 \wedge m)^2$: an $n$-independent floor
    determined by the contamination mass and level;
  \item $n^{-2\beta_1/(2\beta_1+1)}\,
    \varepsilon^{2/(2\beta_1+1)}$: an interaction term coupling
    sample size to contamination structure.
\end{enumerate}
Term~(2) is the key: an irreducible estimation floor that no amount
of data can overcome, proportional to the contamination mass.  In
the Yosida--Hewitt reading, this is the cost of the PFA component:
mass that is invisible to every finite refinement but creates a
permanent bias in any estimator.

\medskip\noindent
\textbf{Plateau detection.}  The transition between the
$n$-dependent regime (term~1) and the $\varepsilon$-dependent floor
(term~2) is detectable from the data.  Lepski's
method~\cite{Lepski1991} and its descendants (Goldenshluger--Lepski)
perform adaptive bandwidth selection that halts when increasing
resolution stops improving the estimate --- i.e., when the estimator
hits a plateau.  The cause of the plateau (contamination mass,
resolution limit, or model misspecification) is not identifiable
from finite data; but the plateau itself is.

This is the precise operational sense in which the finite/countable
boundary is ``felt'' without being ``seen'': the convergence rate of
an estimator encodes which side of the boundary governs the
dominant error term.  If $\sigma$-additive, the rate improves with
$n$; if the PFA component dominates, the rate plateaus.

\medskip\noindent
\textbf{Intrinsic dimension.}  In the classical rate
$n^{-2s/(2s+d)}$, the exponent $d$ is traditionally the ambient
dimension.  The intrinsic-dimension literature replaces $d$ with an
effective dimension adapted to the data geometry: Kpotufe~\cite{Kpotufe2011} shows that $k$-NN
regression achieves $n^{-2s/(2s+\rho)}$ with $\rho$ the local
intrinsic dimension; Yang--Barron~\cite{YangBarron1999}
characterize minimax rates via metric entropy, which encodes
effective dimension.  The programme's $D_\Lambda$ (observational
resolution dimension, the ratio of distinguishability growth to
valuation growth) is the same substitution in the language of valued
refinements.  This is not new mathematics; it is a translation
confirming that the rate-theoretic content of $D_\Lambda$ is the
standard metric-entropy exponent.

\medskip\noindent
\textbf{The recursive undermining.}  One might ask whether a
higher-order information-theoretic criterion could \emph{detect} the
boundary (not just feel its consequences).  The answer is no, by
essentially the same compactness argument that operates at first
order.  Any criterion weaker than $\sigma$-additivity is satisfied
by PFA charges (they pass all finite tests); any criterion strong
enough to exclude PFA charges quantifies over all countable
sequences and is equivalent to $\sigma$-additivity.  There is no
intermediate logic between first-order (compactness, can't see the
boundary) and full second-order (no compactness, sees it by brute
force) that both has compactness and can express $\sigma$-additivity
(a consequence of Lindstr\"{o}m's theorem applied to probability
structures).

\medskip\noindent
\textbf{The Kolmogorov extension mechanism.}  One might hope that
the structural constraint of finite additivity---the PFA component
must marginalize consistently across resolutions---would give a
multi-resolution estimator leverage unavailable against arbitrary
contamination at \emph{fixed} resolution.  It does not.  At any
fixed finite collection of resolution levels $G_1,\ldots,G_K$, the
restrictions of a PFA charge form a consistent projective family
indistinguishable from the marginals of some $\sigma$-additive measure
(by Kolmogorov extension on finite systems).  The contamination
floor $\varepsilon^2$ is therefore the same whether the contamination
is genuinely PFA or arbitrary, at any single resolution snapshot.

However, this indistinguishability is \emph{pointwise}, not
\emph{asymptotic}.  A PFA charge on $B = \bigcup_k G_k$ cannot have
$\sigma$-additive marginals at all levels simultaneously---if it
did, it would agree with the $\sigma$-additive extension on all of $B$,
contradicting pure finite additivity.  In the estimation setting
where $k^*(n) \to \infty$ as $n \to \infty$, the PFA component's
singularity (Yosida--Hewitt) eventually forces its marginals outside
any fixed smoothness class, making the contamination fraction
$\varepsilon$ consistently estimable in the limit~\cite{PatraSen2016, ChenGaoRen2016}.  The boundary is invisible
at each finite stage but visible in the \emph{trajectory} of
increasing resolution---the precise sense in which it is ``felt''
operationally.

\medskip\noindent
\textbf{Synthesis.}  The estimation-theoretic perspective does not
resolve the finite/countable boundary; it translates the boundary
into a rate-theoretic statement.  The translation connects three
literatures that do not normally cite each other:
\begin{itemize}[nosep]
  \item \emph{Probability foundations} (de~Finetti, Howson,
    Seidenfeld): $\sigma$-additivity is not forced by coherence;
    the boundary is philosophically open.
  \item \emph{Robust estimation} (Huber, Liu--Gao, Lepski):
    contamination creates an $n$-independent floor; adaptive methods
    detect the plateau.
  \item \emph{Model theory} (Lindstr\"{o}m, \L{}o\'{s}, Hoover):
    compactness prevents first-order detection; the boundary is
    logically inaccessible.
\end{itemize}
The convergence is that all three traditions agree on the same
structure: the boundary cannot be detected from finite data, but the
\emph{rate} at which finite data converges encodes which side of the
boundary dominates.  Under conditions where the PFA component is
small relative to the observation resolution --- the operational
content of countable additivity in the programme's framework --- the
distinction between ``is $\sigma$-additive'' and ``looks
$\sigma$-additive at every accessible scale'' is vacuous.

% ==================================================================
\section{Stone duality and realization}
\label{sec:realization}
% ==================================================================

\subsection{Classical framework}

\textbf{Stone representation.}  Every Boolean algebra $B$ is
isomorphic to the clopen algebra of a compact, Hausdorff, totally
disconnected space $\St(B)$ (its Stone
space)~\cite{Johnstone1982}.

\textbf{Charge--measure correspondence.}  Every finitely additive
charge on $B$ corresponds to a regular Borel (Baire) measure on
$\St(B)$~\cite[\S326]{Fremlin2003}, \cite[\S53]{Halmos1950}.

\textbf{Support characterization.}  A charge on $B$ is
$\sigma$-additive if and only if the corresponding Stone-space
measure vanishes on meagre Baire sets, equivalently assigns full
mass to countably complete
ultrafilters~\cite[Theorem~10.5.3]{RaoRao1983}.

\textbf{Principal ultrafilters.}  Every principal ultrafilter is
countably complete.  The converse holds for $\mathcal{P}(\kappa)$
when $\kappa$ is below the first measurable cardinal (Ulam's
theorem).  For non-$\sigma$-complete algebras (typical cylinder
algebras), countably-complete non-principal ultrafilters can exist
without large cardinals.

\subsection{The realization problem}

The Stone construction gives an \emph{unconditional} object: a Baire
probability measure $\hat{\mu}$ on $\St(\mathcal{C})$ always exists
from compatible charges, by Stone representation combined with
Choksi's inverse limit theorem.

Descent to a state space $\Omega$ requires that $\hat{\mu}$ be
concentrated on principal ultrafilters:
$\hat{\mu}(\pure(\Omega)) = 1$, where
$\pure(\omega) = \{E \in \mathcal{C} : \omega \in E\}$.  Under
discriminability and the assumption that countably-complete
ultrafilters are principal (which holds for $\sigma$-complete
algebras below the first measurable cardinal), concentration on
principal ultrafilters is equivalent to $\sigma$-additivity of the
component charges.

The three-level nesting captures the structure:
\[
  \pure(\Omega) \;\subseteq\;
  \{\text{countably complete ultrafilters}\} \;\subseteq\;
  \St(\mathcal{C}).
\]
$\sigma$-additivity gives mass on the middle set (intrinsic to the
algebra).  Descent to $\Omega$ requires mass on the inner set
(depends on the choice of realization).

\subsection{The philosophical dimension}
\label{sec:ea-pr-vdr}

Which ultrafilters count as ``realized'' is additional structure
over the algebra.  In the Boolean case, the distinction between the
Stone measure (unconditional) and its descent to $\Omega$ parallels
van Fraassen's~\cite{vanFraassen1980} distinction between empirical
adequacy and realist commitment: the Stone measure is the
``empirically adequate'' object; descent requires choosing which
points are real.  The identification of Stone-space construction
with empirical adequacy (EA), descent with probabilistic realism
(PR), and concentration on dispersion-free states with value-definite
realism (VDR) has no precedent in the literature.  It is a
philosophical contribution connecting van Fraassen to quantum logic
via McDonald--Bimb\'{o} duality, not a mathematical theorem.

% ==================================================================
\section{Value-definiteness and quantum logic}
\label{sec:vdr}
% ==================================================================

\subsection{Distributivity and ultrafilters}

Boolean algebras admit 2-valued homomorphisms (ultrafilters) by
Zorn's lemma.  These are the dispersion-free states: they assign
definite values (0 or 1) to every proposition simultaneously.

The Kochen--Specker theorem~\cite{KochenSpecker1967} shows that
for the lattice $L(H)$ of closed subspaces of a Hilbert space $H$
with $\dim H \geq 3$, no dispersion-free state exists.
Non-distributivity is the structural cause: $L(H)$ is orthomodular
but not distributive, and the failure of distributivity blocks the
lattice-theoretic argument that guarantees ultrafilters in the
Boolean case.

Gleason's theorem~\cite{Gleason1957} completes the picture for
$L(H)$, $\dim H \geq 3$: every $\sigma$-additive state is of the
form $s(P) = \mathrm{tr}(\rho P)$ for a unique density operator
$\rho$.  (The theorem fails for $\dim H = 2$, where
dispersion-free states on $L(\mathbb{C}^2)$ do exist.)
Value-definiteness is blocked in dimension $\geq 3$, but
probabilistic realism (states as density operators) is available.

The Bub--Clifton theorem~\cite{BubClifton1996},
\cite{HalvorsonClifton1999} identifies the maximal Boolean
subalgebra $D(\rho, R) \subseteq L(H)$ on which a given state is a
mixture of dispersion-free states.  This gives the precise boundary
of value-definiteness within a non-Boolean algebra.

\subsection{Orthomodular lattice duality}

McDonald--Bimb\'{o}~\cite{McDonaldBimbo2023} establish a topological
duality for orthomodular lattices analogous to Stone duality for
Boolean algebras.  A key structural difference: principal filters
$P(A)$ are part of the dual structure in the OML case (unlike Stone
duality where they are invisible).  This means realization is
\emph{constrained} structure in the OML case, not \emph{additional}
structure as in the Boolean case.  As of 2026, this duality has received essentially no attention in
quantum foundations.

Gunji--Haruna--Uragami~\cite{Gunji2026} show that gluing Boolean
algebras via categorical pushout produces orthomodular lattices,
generically non-distributive.

\subsection{Topos approaches}

Three programmes reformulate quantum logic in topos-theoretic terms:

\begin{itemize}
  \item \textbf{Isham--Butterfield~\cite{IshamButterfield1998}:}
    the topos $\mathbf{Sets}^{\mathcal{V}(\mathcal{H})^{op}}$ of
    presheaves on the poset of commutative subalgebras.  The spectral
    presheaf $\underline{\Sigma}$ assigns to each Boolean context its
    Gelfand spectrum; it has no global elements $=$ Kochen--Specker.

  \item \textbf{D\"{o}ring--Isham~\cite{DoringIsham2008}:}
    Daseinisation maps a projection to a sub-object of
    $\underline{\Sigma}$ by approximating within each Boolean context.
    This goes \emph{up} (approximating non-Boolean by Boolean);
    descent goes \emph{down} (Stone space to realized points).

  \item \textbf{Heunen--Landsman--Spitters~\cite{HeunenLandsmanSpitters2009}:}
    Bohrification encodes the same structure via spectral presheaves.
    More powerful than the single-dual-space approach (presheaf over
    contexts, not single Stone space) but connects less directly to
    van Fraassen's framework.
\end{itemize}

R\'{e}dei--Summers~\cite{RedeiSummers2006} review quantum
probability in the von Neumann algebra framework: classical
probability $=$ abelian von Neumann algebra; quantum $=$ nonabelian;
normal states $=$ $\sigma$-additive.

\subsection{The OML extension problem}

\subsubsection{State space universality}

Shultz~\cite{Shultz1974} proves that orthomodularity imposes
\emph{no} geometric constraint on state spaces: every compact convex
subset of a locally convex space arises as the state space of some
$\sigma$-orthomodular lattice.  The construction uses Greechie
pasting ($G$-spaces): points and frames with loop constraints.
However, $\sigma$-additive state spaces can be geometrically
degenerate --- the Lebesgue measure algebra has a $\sigma$-state
space with no extreme points, hence not affinely equivalent to any
compact convex set~\cite[p.~327]{Shultz1974}.

\subsubsection{Exotic logics and the Greechie pasting technique}

Pt\'{a}k~\cite{Ptak1987} constructs \emph{exotic logics}:
$\sigma$-orthocomplete OMPs admitting finitely additive states but
no $\sigma$-additive states.  The construction uses Greechie
pasting: atom identifications across Boolean blocks create
contradictory partition-of-unity constraints at the
$\sigma$-additive level.  Each block individually admits
$\sigma$-additive states; the global pasting forbids them.

Navara~\cite{Navara1992} shows that regularity ($P$-regularity,
inner regularity with respect to finite-dimensional projections)
does \emph{not} force $\sigma$-additivity on general OMPs.  He
modifies Pt\'{a}k's construction to produce exotic logics on which
all states are $P$-regular yet none are $\sigma$-additive.  (The
B\'{e}aver--Cook result that regularity forces $\sigma$-additivity
is special to von Neumann algebras.)

\medskip\noindent
\textbf{Synthesis.}  Greechie pasting is the universal construction
tool for non-Boolean phenomena:
Pt\'{a}k~\cite{Ptak1987} uses it to produce exotic logics,
Shultz~\cite{Shultz1974} to prove state-space universality,
Navara~\cite{Navara1992} to kill regularity as an admissibility
condition, and Wilce~\cite{Wilce2009} surveys the technique
systematically.  The obstruction to $\sigma$-additive extension in
the non-Boolean case is combinatorial/geometric (from the pasting),
structurally different from the Boolean case where the obstruction
is analytic (weak $(\sigma,\infty)$-distributivity).

\subsubsection{The structural obstruction}

Pt\'{a}k--Pulmannov\'{a}~\cite{PtakPulmannova1994} prove: if an
orthomodular lattice possesses a unital set of subadditive probability
measures, it must be Boolean (distributive).  Conditions rich enough
to force $\sigma$-additivity destroy the non-Boolean structure.

Related results exist under strong completeness conditions:
Bunce--Wright~\cite{BunceWright1992} solve the Mackey--Gleason problem
for von Neumann algebras (every bounded finitely additive measure on
the projection lattice extends to a bounded linear functional),
and Chetcuti--Dvure\v{c}enskij~\cite{ChetcutiDvurecenskij2003} use
finitely additive states to characterize completeness of inner product
spaces.  Neither directly addresses the FA-to-$\sigma$-additive
extension question, and no result bridges the gap between ``too weak''
(allowing non-$\sigma$-additive states) and ``too strong'' (forcing
Booleanity).

\begin{openproblem}[OML extension]
\label{op:oml-extension}
Let $A$ be an orthomodular lattice and
$s \colon A \to [0,1]$ a state.  Let $S_0(A)$ be the
McDonald--Bimb\'{o} dual space with clopen $\perp$-stable sets
$\mathrm{CO}(S_0(A))^\dagger$, and define $\mu(h(a)) = s(a)$ where
$h(a) = \{x \in F(A) : a \in x\}$.
\begin{enumerate}[label=(\alph*),nosep]
  \item Under what conditions on $A$ does $\mu$ extend to a
    $\sigma$-additive measure on the Baire $\sigma$-algebra of
    $S_0(A)$?
  \item Does any condition strictly weaker than Booleanity suffice?
\end{enumerate}
\end{openproblem}

For $A = L(H)$ with $\dim H \geq 3$, Gleason gives the extension.
For general OMLs, Gleason-type results are unavailable and
Carath\'{e}odory does not apply (the clopens form an OML, not a
Boolean algebra).

% ==================================================================
\section{Delay reconstruction}
\label{sec:reconstruction}
% ==================================================================

\subsection{Embedding theorems}

\textbf{Takens (1981).}  For a smooth compact manifold $M$ with
$C^2$ diffeomorphism $T$ and generic $C^2$ observable $h$, the delay
map $\Phi_L(x) = (h(x), h(Tx), \ldots, h(T^{L-1}x))$ is an
injective immersion for
$L \geq 2\dim M + 1$~\cite{Takens1981}.

\textbf{Sauer--Yorke--Casdagli (1991).}  Prevalence-based
generalization: the set of observables for which $\Phi_L$ is an
embedding has full prevalence measure for $L > 2d_B(X)$ where
$d_B$ is the box-counting dimension~\cite{SauerYorkeCasdagli1991}.

\textbf{Botvinick-Greenhouse et al.\ (2025).}  Measure-theoretic
Takens theorem via optimal transport: if the pushforward measures
converge in Wasserstein distance, the delay map is a
measure-theoretic embedding~\cite{BotvinickGreenhouse2025}.

\subsection{Generating partitions and algebraic reconstruction}

The ergodic-theoretic parallel: an observable $\sigma$-algebra
equals the full $\sigma$-algebra mod $\mu$ if and only if the delay
algebra is dense in $L^2$, if and only if the delay map is a
measure-theoretic embedding.  This is Krieger's generating partition
theorem~\cite{Krieger1970} in delay-map language
(cf.~Sinai~\cite{Sinai1959} for the entropy definition,
Rokhlin~\cite{Rokhlin1967} for the metric).

The Rokhlin distance $\delta(\mathcal{O})$ measures how close the
observable $\sigma$-algebra is to generating everything.
Collision probability ($=$R\'{e}nyi-2 entropy of the fibre
decomposition) measures the resolution structure.

\textbf{Geometric $\neq$ algebraic reconstruction.}
Collision convergence (geometric) and $\delta \to 0$ (algebraic)
are genuinely inequivalent.  A skew-product counterexample
($X = A^{\mathbb{Z}} \times B^{\mathbb{Z}}$, $T = \sigma \times
\sigma$, $h(a,b) = a_0$) gives collision $= 2^{-(L+1)} \to 0$
(fibres shrink in the $a$-coordinate) but $\delta_L = 1/4$ for all
$L$ ($\{b_0 = 0\}$ is maximally unresolved).  This is the one
observation from this programme that I have not found in the prior
literature, though it may well be known to ergodic theorists.

\subsection{Practical diagnostics}

\begin{itemize}[nosep]
  \item \textbf{False nearest neighbours (FNN):}
    embedding dimension selection~\cite{KennelBrownAbarbanel1992},
    sensitive to noise~\cite{RhodesMorari1997}.

  \item \textbf{Ragwitz--Kantz:} local prediction error minimization
    for $(d, L)$ selection~\cite{RagwitzKantz2002}.

  \item \textbf{Billings--Voon:} residual autocorrelation as
    model validity test~\cite{BillingsVoon1986}.

  \item \textbf{Casdagli:} deterministic vs.\ stochastic modelling
    discrimination~\cite{Casdagli1992}.
\end{itemize}

Every practical criterion implicitly assumes a modelling commitment:
deterministic criteria fail under noise because noise violates the
commitment.  Finite-sample convergence rates are
$n^{-s/(2s+d)}$~\cite{Stone1982}; elbow stopping via
Lepski~\cite{Lepski1991}.

\subsection{The reconstruction problem in programme context}

The delay setting is the primary source of observation algebras for
this programme: a dynamical system with a scalar observable produces
a directed system of Boolean algebras (delay algebras at increasing
embedding dimension).  The three-layer question then becomes: under
what conditions on the dynamics and observable do the resulting
charges extend ($\sigma$-additivity), descend (realization), and
determine definite values (value-definiteness)?

In the Boolean case (classical dynamics), all three are available in
principle; the obstruction is purely at the extension level
($\sigma$-additivity).  In the quantum case (non-commutative
dynamics), value-definiteness is blocked by Kochen--Specker and the
extension problem takes a different form
(Problem~\ref{op:oml-extension}).

% ==================================================================
\section{Ergodic theory of finitely additive charges}
\label{sec:ergodic}
% ==================================================================

The Yosida--Hewitt decomposition commutes with shift-invariance:
if $\ell$ is a $T$-invariant charge, both $\ell_c$ and $\ell_p$ are
$T$-invariant.  The structure of the $\sigma$-additive part is the
subject of classical ergodic theory (Birkhoff, von Neumann).  The
purely finitely additive part is far less understood.

\begin{openproblem}[Extreme PFA shift-invariant charges]
\label{op:ergodic-pfa}
Let $T \colon \Omega \to \Omega$ be a measurable transformation on a
standard Borel space.  What are the extreme points of the convex set
of $T$-invariant purely finitely additive charges on $\Omega$?
\end{openproblem}

This connects to the structure of Banach limits and invariant means.
The problem appears out of reach with current techniques.

% ==================================================================
\section{Summary of the boundary}
\label{sec:summary}
% ==================================================================

The extension boundary takes different forms in different algebraic
settings:

\begin{center}
\begin{tabular}{@{}llll@{}}
\toprule
\textbf{Setting} & \textbf{Obstruction} &
  \textbf{Characterization} & \textbf{Status} \\
\midrule
Boolean, $\sigma$-complete & analytic &
  KVP (391D) & classical \\
Boolean, non-$\sigma$-complete & analytic/topological &
  Problem~\ref{op:strategy-d} & open \\
OML, $L(H)$, $\dim \geq 3$ & algebraic &
  Gleason & classical \\
OML, general & combinatorial &
  Problem~\ref{op:oml-extension} & open \\
\bottomrule
\end{tabular}
\end{center}

The Boolean and non-Boolean extension problems are structurally
independent: the Boolean obstruction is analytic (weak
$(\sigma,\infty)$-distributivity, or existence of measure-free
factors), while the OML obstruction is combinatorial (Greechie
pasting creating contradictory normalization).  Neither subsumes the
other.  In the Boolean non-$\sigma$-complete case, every known ZFC
measure-free construction is either $\sigma$-complete or has atoms
(Table~\ref{tab:near-miss}); the exact Strategy~D parameter regime
is likely ZFC-independent.%
\footnote{\textbf{Superseded 2026-06-18} --- see the footnote at the
``Strategy~D is likely ZFC-independent'' claim earlier: the regime has a positive
ZFC answer (Gaifman~1964) and the Table's ``$\sigma$-complete'' entries read
$\sigma$-completeness off the \emph{completed}/Gleason-cover form, not the base
algebra. Strategy~D retired as a research target.}

Sheaf-theoretic, topos-theoretic, and locale-theoretic approaches
to the Boolean extension problem converge on dictionary translations
rather than theorems: the finite/countable boundary is invisible to
finitary frameworks, and infinitary frameworks restate classical
conditions.  The OML extension problem has received less attention
from these perspectives; the McDonald--Bimb\'{o} duality provides
the natural framework but the extension question within it remains
unasked in the literature.
